% Title: Chili by Wes
% Description: Slightly more intesting tradiontal Midwest chili
% Author: Wesley Jones
% Date: December 2017

\documentclass{article}

\usepackage[nonumber,noindex,nocontents]{cuisine}
\usepackage[margin=1in]{geometry}
\usepackage{endnotes}

% Set widths
\RecipeWidths  
  {1\textwidth}    %{Total recipe width}             
  {0.1\textwidth}  %{Number of servings width}                   
  {0.05\textwidth} %{Step number width}             
  {0.15\textwidth} %{Ingredient width} 
  {0.1\textwidth}  %{Quantity width} 
  {0.1\textwidth}  %{Units width}


% End notes instead of footnotes
\let\footnote=\endnote


\begin{document}

\begin{recipe}{Chili by Wes}{\textit{Serves: 8}}{\textit{Preparation Time: 20 minutes}}
  \freeform
    Chili should be straightforward.
    It should be warm, not hot.
    It should be hearty, not thick.
    It is ubiquitous --- so here's my way.

  \freeform
    \textbf{\textit{Seasoning}}

    \ing[2]{Tbsp}{chili powder}
    \ing[1]{tsp}{cumin}
    \ing[1]{tsp}{cocoa powder}
    \ing[1]{tsp}{red pepper flakes}
    \ing[\fr{1}{2}]{tsp}{garlic powder}
    \ing[\fr{1}{2}]{tsp}{salt}
    \ing[\fr{1}{2}]{tsp}{paprika}
    \ing[\fr{1}{4}]{tsp}{cayenne}
    Combine all in a small dish.
    Set aside.

  \freeform
    \textbf{\textit{Soup}}

    \ing[1]{lb}{ground beef}
    \ing[1]{medium}{onion, diced}
    \ing[1]{large}{green bell pepper, diced}
    \ing[2]{cloves}{garlic, minced}
    \ing[28]{oz}{canned, whole tomatoes}
    \ing[1]{dish}{seasonings from Step 1}
    Brown beef in 4 to 6 quart dutch oven or pot.
    As the beef finishes add pepper, onion, and garlic.
    Allow vegetables to simmer and soften.
    Clear a spot and add the seasonings directly into the drippings and exposed bottom of the pan.
    Allow these to bloom for 1-2 minutes before stirring into mixture.
    Add can tomatoes and cover for 15 minutes.

    \ing[16]{oz}{tomato sauce}
    \ing[16]{oz}{pinto beans, drained \& risned}
    \ing[16]{oz}{chicken, beef, or vegetable broth}
    After tomatoes have simmered and softened, break them down using a wooden spoon or other suitable utensil.
    Break them up as fine as desired.
    Add remaining ingredients, and allow to come to a simmer.
    Cook for longer to allow flavors to ``meld''.

\end{recipe}

\end{document}
